\mxunit{automatismes-ch13}{Consolidation, automatismes et transition vers la Première}

\begin{fluencysheet}{mar.~1\textsuperscript{er}~juin~2027}{Problèmes de synthèse mobilisant plusieurs domaines}
\mxfq{Échauffement · Flash}{EQ}{Résoudre $3(x-1)=x+5$.}{$3x-3=x+5$, $2x=8$, $x=4$.}{}
\mxfq{Cœur · Méthode}{GEO}{Pour quelle valeur de $k$ les points $A\pt{1;1}$, $B\pt{3;5}$ et $C\pt{k;9}$ sont-ils alignés\,?}{$\vect{AB}\vc24$, $\vect{AC}\vc{k-1}{8}$~: $2\times8-4(k-1)=0$, donc $k=5$.}{}
\mxfq{Cœur + · Méthode}{PRO}{Un prix augmente de $\pc{10}$ puis diminue de $\pc{10}$. Quelle est l'évolution globale\,?}{$\num{1,1}\times\num{0,9}=\num{0,99}$~: baisse de $\pc{1}$.}{}
\mxfq{Retour · Flash}{PER}{Quelle est l'aire totale d'un pavé droit de dimensions $4$, $3$ et $2$ (en cm)\,?}{$2(4\times3+4\times2+3\times2)=52$ cm$^2$.}{}
\mxfq{Fil rouge · Flash}{SOL}{Calculer le volume exact d'une boule de rayon $3$ cm.}{$V=\dfrac43\pi\times27=36\pi$ cm$^3$.}{}
\end{fluencysheet}

\begin{fluencysheet}{jeu.~3~juin~2027}{Bilan des six compétences}
\mxfq{Échauffement · Métacognition}{MET}{Citer les six compétences mathématiques travaillées au lycée.}{Chercher, modéliser, représenter, raisonner, calculer, communiquer.}{}
\mxfq{Cœur · Métacognition}{MET}{Associer~: «~traduire un énoncé par une équation~»\,; «~tracer une courbe~»\,; «~justifier une égalité par un calcul rédigé~» à une compétence.}{Modéliser\,; représenter\,; raisonner (ou communiquer, selon l'accent mis sur la rédaction).}{}
\mxfq{Cœur + · Métacognition}{MET}{Pour résoudre $2x+3=11$, quelle compétence est principalement mobilisée\,?}{Calculer (manipuler une équation du premier degré).}{À discuter oralement~: la mise en équation d'un problème relève de «~modéliser~».}
\mxfq{Retour · Problème}{MOD}{De l'eau à $20$ °C chauffe de $3$ °C par minute~: $T(t)=20+3t$. Au bout de combien de minutes atteint-elle $65$ °C\,?}{$20+3t=65\Leftrightarrow t=15$ min.}{}
\mxfq{Fil rouge · Flash}{EQ}{Résoudre dans $\R$~: $\num{0,5}x+\num{1,2}=\num{2,7}$.}{$x=3$}{}
\end{fluencysheet}

\begin{fluencysheet}{lun.~7~juin~2027}{Réactivation ciblée des prérequis pour la Première}
\mxfq{Échauffement · Flash}{DEV}{Développer $(x+3)^2$.}{$x^2+6x+9$.}{}
\mxfq{Cœur · Méthode}{EQ}{Vérifier que $x^2-5x+6=(x-2)(x-3)$ et en déduire les solutions de $x^2-5x+6=0$.}{$(x-2)(x-3)=x^2-5x+6$\,; produit nul~: $x=2$ ou $x=3$.}{}
\mxfq{Cœur + · Anticipation}{MOD}{Pour $f(x)=x^2$, calculer le taux d'accroissement $\dfrac{f(3)-f(1)}{3-1}$.}{$\dfrac{9-1}{2}=4$~: prépare la notion de nombre dérivé en Première.}{Notion de Première, simple calcul de réinvestissement.}
\mxfq{Retour · Programme de construction}{CON}{On place $M$ et $N$ milieux de $[AB]$ et $[AC]$ puis on trace $[MN]$. Que peut-on dire de $(MN)$ et $(BC)$\,? Calculer $MN$ si $BC=9$ cm.}{$(MN)\parallel(BC)$ et $MN=\dfrac92=\num{4,5}$ cm.}{}
\mxfq{Fil rouge · Rédiger}{DEM}{Démontrer qu'un triangle de côtés $7$, $24$ et $25$ est rectangle.}{$7^2+24^2=49+576=625=25^2$~: réciproque de Pythagore.}{}
\end{fluencysheet}

\begin{fluencysheet}{mar.~8~juin~2027}{Problèmes de synthèse~: choisir et justifier une stratégie}
\mxfq{Échauffement · Lecture de code}{ALG}{Que renvoie $f(2)$\,?\mxpseudo{def f(x):\\ \mxtab return x**2 - 3*x + 1}}{$4-6+1=-1$.}{}
\mxfq{Cœur · Complétion}{ALG}{Compléter pour que la fonction renvoie \code{True} lorsque $A$, $B$, $C$ sont alignés~:\mxpseudo{def alignes(xa, ya, xb, yb, xc, yc):\\ \mxtab return (xb-xa)*(yc-ya) - (yb-ya)*(xc-xa) == ....}}{On compare le déterminant à $0$~: \code{== 0}.}{}
\mxfq{Cœur + · Complétion}{ALG}{Compléter la fonction valeur absolue~:\mxpseudo{def vabs(x):\\ \mxtab if x \textgreater{}= 0:\\ \mxtab\mxtab return x\\ \mxtab else:\\ \mxtab\mxtab return ....}}{\code{return -x}.}{Contrôle~: \code{vabs(-4)} renvoie $4$.}
\mxfq{Retour · Flash}{NUM}{Écrire en notation décimale $\num{2,5}\times10^3$ puis $\num{6,2}\times10^{-2}$.}{$\num{2500}$\,; $\num{0,062}$.}{}
\mxfq{Fil rouge · Flash}{REP}{Donner le symétrique de $A\pt{2;-3}$ par rapport à l'origine.}{$\pt{-2;3}$.}{}
\end{fluencysheet}

\begin{fluencysheet}{jeu.~10~juin~2027}{Présentation des enseignements de spécialité\,; bilan}
\mxfq{Échauffement · Flash}{NUM}{Exprimer $\dfrac38$ en pourcentage.}{$\pc{37,5}$.}{}
\mxfq{Cœur · Métacognition}{MET}{Quelles sont tes trois priorités de révision avant la Première\,?}{Réponse personnelle~: repérer dans les fiches de l'année les questions où l'on a le plus hésité.}{}
\mxfq{Cœur + · Flash}{EQ}{Résoudre $2x^2=18$.}{$x^2=9$~: $x=3$ ou $x=-3$.}{}
\mxfq{Retour · Piège}{PUI}{Calculer $(-2)^4$ et $-2^4$.}{$16$ et $-16$.}{Les parenthèses décident à qui s'applique l'exposant.}
\mxfq{Fil rouge · Vrai/faux}{DIV}{Vrai ou faux\,? Justifier. Si un entier est divisible par $4$ et par $6$, alors il est divisible par $24$.}{Faux. Contre-exemple~: $12$ est divisible par $4$ et par $6$, pas par $24$\,; il est divisible par $\ppcm\pt{4;6}=12$.}{}
\end{fluencysheet}
